% Auto-generated by generate_tex_fragments.py -- do not hand-edit.
\begin{longtable}{@{}L{2.0cm}L{2.0cm}L{2.1cm}L{6.2cm}L{1.9cm}@{}}
\caption{P7: Repairability is necessary because relevant information emerges over time. Record, design, population/setting, finding as charted (verbatim from the library record), and direction relative to the proposition. n=6 linked records (kg/graph.json edges).}\label{tab:p7}\\
\toprule
\textbf{Record} & \textbf{Design} & \textbf{Population\slash{}setting} & \textbf{Finding (as charted)} & \textbf{Direction}\\
\midrule
\endfirsthead
\multicolumn{5}{l}{\textit{P7 continued}}\\
\toprule
\textbf{Record} & \textbf{Design} & \textbf{Population\slash{}setting} & \textbf{Finding (as charted)} & \textbf{Direction}\\
\midrule
\endhead
\bottomrule
\endfoot
Michau L 2021 \cite{ref113} & Programme-theory\slash{}programme-\seqsplit{description} paper (non-empirical evaluation article) & Raising Voices programme learning, Uganda-originated model adapted for broader use & The paper describes the rationale for revising SASA! after ten years of programme learning and research, detailing how the revised 'SASA! Together' model \seqsplit{incorporated} field learning, and reflects on debates such as how best to scale up violence-prevention programmes. & Supports\\
\\[3pt]
Silliman B 1999 \cite{ref116} & Narrative review & General review of marriage-\seqsplit{preparation} program practice, US & Reviews two decades of research on marriage-\seqsplit{preparation} programs and gives design \seqsplit{recommendations} covering needs assessment, timing\slash{}duration, format, content, provider \seqsplit{characteristics}, and promotion; does not report new outcome data. & Supports\\
\\[3pt]
Subramanee SD 2022 \cite{srp36429857} & Systematic review (PRISMA, 13 included studies of 520 screened), PROSPERO \seqsplit{CRD42020190410} & South Asian \seqsplit{populations} (India, Bangladesh, Pakistan, Nepal, Sri Lanka), including Muslim subgroups & Consistent correlates of child marriage were rural residence, low education, poor economic background, low mass-media exposure and religion (Hindu and Muslim in particular countries); downstream maternal-health effects included lower antenatal care use, \seqsplit{institutional} delivery and skilled birth attendance. & Supports\\
\\[3pt]
Upadhyay UD 2014 \cite{ref16} & Instrument \seqsplit{development} and validation (cross-sectional survey, factor analysis) & 1,892 women at 13 family planning and 6 abortion facilities, USA & Generated 26 candidate items, retained 14 items across 3 subscales (freedom from coercion, \seqsplit{communication}, decision-making) via factor analysis; the freedom-from-coercion and \seqsplit{communication} subscales were inversely associated with \seqsplit{unprotected} sex in the past 3 months, supporting construct validity. & Supports\\
\\[3pt]
Clyde TL 2020 \cite{ref115} & Narrative review\slash{}position paper & General discussion of \seqsplit{contemporary} premarital-education practice, US-centric & Reviews social trends (\seqsplit{individualism}\slash{}commitment \seqsplit{ambivalence}, changing marriage attitudes, premarital \seqsplit{relationship} histories, media \seqsplit{environment}) and argues current premarital \seqsplit{interventions} need updated content and structure to remain relevant to engaged couples today; offers four proposals for \seqsplit{redesigning} premarital \seqsplit{interventions}. & Relevant only\\
\\[3pt]
Upadhyay UD 2021 \cite{ref18} & Instrument \seqsplit{development} and validation (\seqsplit{qualitative} item generation + national survey, EFA, logistic regression) & 1,117 sexually active young people aged 15-24, USA (national sample) & Developed a 23-item, 7-subscale Sexual and \seqsplit{Reproductive} \seqsplit{Empowerment} Scale for \seqsplit{Adolescents}\slash{}Young Adults (including 'choice of partners, marriage, and children' as a named subscale); validation via logistic regression supported construct validity. & Relevant only\\
\\[3pt]
\end{longtable}
